\chapter{Introducción}
\label{cap:introduccion}

\section{Problema. Motivación}
\label{sec:intro-motivacion}

Los agentes basados en grandes modelos de lenguaje ya toman decisiones de inversión. Con ese despliegue ha crecido una demanda concreta: supervisar de forma continua y trazable lo que la IA agéntica decide. A medida que estos sistemas ganan autonomía, sus decisiones se vuelven difíciles de seguir para un humano, y crece el interés en los \emph{guardrails}: una supervisión automática que acote lo que el agente puede hacer.

Muchas personas ya confían su dinero a inversiones gestionadas por estos agentes. Un agente LLM no es determinista: ante la misma pregunta puede responder de forma distinta, y no señala cuándo acierta o falla. En la gestión de capital, donde un solo error tiene coste elevado, esa indeterminación es crítica.

El sector financiero exige por ello un control más estricto. Los agentes LLM ejecutan las instrucciones con fidelidad antes que maximizar el beneficio \cite{lopezlira2025cantrade}, su confianza declarada está mal calibrada y son manipulables. Entre la supervisión que se reclama y una herramienta falsable que la ofrezca hay una distancia que este trabajo aborda.


El agente estudiado es una réplica de AI Hedge Fund: cinco personalidades inversoras sobre un mismo modelo de lenguaje que deciden, para cada activo y cada día, una dirección y un tamaño. Sin supervisión pierde dinero. Sobre SPY acierta la dirección el $0{,}366$ de los días, por debajo de lo que daría una moneda, y una inversión inicial de \euro{}1000 termina en \euro{}700.\,% fuente: cap. STRATA-SPY (SPY M5: acc 0,366, equity 0,70)
El fenómeno no es propio de SPY: el agente acierta por debajo del medio punto en siete de los diez activos del panel.\,% fuente: cap. panel (7/10 < 0,5)

La pérdida tiene estructura. El agente mantiene una posición casi rígida que apenas atiende a la señal del día. Un fallo sistemático se puede corregir desde fuera; ese es el problema que aborda este trabajo.

\section{Propuesta: supervisión estadística externa}
\label{sec:intro-idea}

STRATA es la solución a este problema: un supervisor estadístico exógeno que audita cada decisión en tiempo de ejecución y la rectifica cuando contradice el comportamiento del mercado. Se apoya en tres detectores ortogonales, cada uno sobre un modelo estadístico clásico que el capítulo~\ref{cap:strata} formaliza: RAM mira la coherencia de la apuesta con el régimen de mercado, PSA la coherencia temporal de las decisiones del agente y GSO la compatibilidad del tamaño con la volatilidad. Cuando uno se dispara, STRATA actúa sobre la decisión: la registra, atenúa su tamaño o la sustituye. La corrección se contrasta de forma pareada contra el mismo agente sin supervisar, decisión a decisión.

La idea de una etapa secundaria que vigila a un modelo primario no es nueva: el \emph{meta-labeling} de López de Prado \cite{lopezdeprado2018} añade un segundo modelo que decide si actuar sobre la señal del primero. STRATA recoge esa lógica y la lleva más lejos. Aquí el primario es el agente LLM, y el supervisor son tres detectores estadísticos clásicos que juzgan su decisión desde ejes independientes y la registran, atenúan o sustituyen. La intervención es exógena: depende de lo que los detectores leen del mercado, no de lo que el agente declare. Eso la separa del equipo de riesgo de TradingAgents \cite{xiao2024tradingagents}, un agente deliberativo interno que comparte la información y los sesgos del sistema.
